% =============================================================================
% ch08_results.tex — Chapter 8: Multi-Sensor Fusion: Camera-LiDAR-IMU-GNSS | NavigatorCrew
% Compiler: XeLaTeX | Language: Hebrew primary, English technical terms via \en{}
% =============================================================================

\selectlanguage{hebrew}

\section{\en{Fusion} רב-חיישני: מצלמה-\en{LiDAR}-\en{IMU}-\en{GNSS}}
\label{sec:results}

% ── 8.1 Tightly-Coupled Fusion ──────────────────────────────────────────────
\subsection{עקרונות \en{Fusion} הדוק (\en{Tightly-Coupled})}
\label{subsec:tightly_coupled}

\en{Fusion} רב-חיישני ל-\en{SLAM} הפך לחיוני להשגת אומדן מצב חזק
ונטול סחיפה בסביבות מאתגרות שבהן כל מודאליות חיישן בודדת נכשלת.
הפרדיגמה הדומיננטית היא \en{Fusion} הדוק (\en{Tightly-Coupled Fusion}),
שבו מדידות גולמיות מכל החיישנים מאופטימזות במשותף בפונקציית עלות
אחת, במקום חיבור רופף של אומדנים בלתי תלויים.

\en{LIO-SAM} מדגים גישה זו על ידי שילוב \en{LiDAR}, \en{IMU}
ו-\en{GPS} במסגרת \en{Factor Graph}. ה-\en{IMU} מספק אילוצי
תנועה בתדר גבוה (200-400 הרץ) המבטלים עיוותים בסריקות \en{LiDAR}
ומאכפים את האופטימיזציה בין \en{Keyframes} של \en{LiDAR}.
\en{GPS} מספק אילוצי מיקום מוחלטים המבטלים סחיפה ארוכת טווח.
\en{LVI-SAM} מרחיב גישה זו לכלול אודומטרייה חזותית-אינרציאלית
כתהליך אומדן מקבילי, היוצר מערכת שיכולה לפעול גם בסביבות עשירות
ב-\en{LiDAR} וגם בסביבות עשירות במידע חזותי.

\begin{figure}[H]
    \centering
    \includegraphics[width=0.92\columnwidth]{figures/fig_lio_sam_factor_graph.png}
    \caption{גרף הגורמים של \en{LIO-SAM}: צמתי \en{IMU} (ירוק),
             צמתי \en{LiDAR} (כחול), צמתי \en{GPS} (אדום),
             וסגירת לולאה (סגול).}
    \label{fig:lio_sam_factor_graph}
\end{figure}

% ── 8.2 R2LIVE: Iterated Error-State Kalman Filter ──────────────────────────
\subsection{\en{R2LIVE}: פילטר קלמן איטרטיבי מבוסס \en{Error-State}}
\label{subsec:r2live}

\en{R2LIVE} מציג פילטר קלמן איטרטיבי הדוק
(\en{Iterated Error-State Kalman Filter}) הממזג מדידות \en{LiDAR},
חזותיות ואינרציאליות ברמת המדידה הגולמית. החדשנות המרכזית היא
שימוש ב-\en{IESKF} המבצע ליניאריזציה סביב האומדן הטוב ביותר
הנוכחי, ובכך נמנע מבעיות העקביות של \en{Fusion} מבוסס
\en{EKF} סטנדרטי.

\en{R2LIVE} משיג \en{ATE} של 0.08 מטר על מערך הנתונים
\en{M2DGR}, תוצאה טובה יותר מ-\en{LIO-SAM} (0.18 מטר)
ומ-\en{LVI-SAM} (0.12 מטר). האלגוריתם פועל בתדר של 20 הרץ
על \en{NVIDIA Jetson AGX Xavier}, מה שהופך אותו למתאים
ליישומי רחפן בזמן אמת. עם זאת, \en{R2LIVE} דורש כיול מדויק
במיוחד בין כל החיישנים, והביצועים מתדרדרים משמעותית כאשר
איכות הכיול ירודה.

\begin{equation}
    \hat{\mathbf{x}}_{k+1} = \hat{\mathbf{x}}_k \oplus \delta\mathbf{x},
    \quad \delta\mathbf{x} \sim \mathcal{N}(\mathbf{0}, \mathbf{P}_k)
    \label{eq:ieskf_state}
\end{equation}

משוואה~\ref{eq:ieskf_state} מגדירה את ייצוג המצב ב-\en{IESKF},
כאשר $\oplus$ היא פעולת החיבור בחבורת \en{Lie} ($SE(3)$),
ו-$\delta\mathbf{x}$ היא השגיאה האינווריאנטית.

\begin{equation}
    \mathbf{K} = \mathbf{P} \mathbf{H}^\top (\mathbf{H} \mathbf{P} \mathbf{H}^\top + \mathbf{R})^{-1},
    \quad \delta\mathbf{x}^+ = \mathbf{K} (\mathbf{z} - h(\hat{\mathbf{x}}))
    \label{eq:ieskf_update}
\end{equation}

משוואה~\ref{eq:ieskf_update} מציגה את עדכון ה-\en{IESKF},
הדומה במבנהו ל-\en{EKF} הסטנדרטי אך פועל על מרחב השגיאה
האינווריאנטית במקום על מרחב המצב עצמו.

% ── 8.3 Challenges and Limitations ──────────────────────────────────────────
\subsection{אתגרים ומגבלות של \en{Fusion} רב-חיישני}
\label{subsec:fusion_challenges}

כשל מרכזי של מערכות \en{Fusion} רב-חיישני הוא חוסר יישור חיישנים
(\en{Sensor Misalignment}), הן מבחינה מרחבית והן מבחינה זמנית.
שגיאות כיול קטנות של 0.1 מעלות בסיבוב או 1 ס״מ בתרגום עלולות
לגרום לסחיפה משמעותית לאורך מסלולים ארוכים. כיול זמני מדויק
חיוני במיוחד: הפרש זמנים של 10 אלפיות שנייה בין מדידת \en{LiDAR}
למדידת \en{IMU} עלול לגרום לשגיאות מיקום של עשרות סנטימטרים
בתנועה מהירה.

אתגר נוסף הוא נשירת חיישנים (\en{Sensor Dropout}). כאשר \en{GPS}
אובד בקניונים עירוניים, \en{LiDAR} נכשל בערפל או מצלמה מאבדת
מעקב בתנאי תאורה נמוכה, המערכת חייבת להתדרדר באופן הדרגתי
ולא להתבדר. גישה מקובלת היא שימוש ב-\en{Weighted Fusion},
שבו משקל כל חיישן מותאם דינמית בהתאם לאיכות המדידה הנוכחית.
גישה נוספת היא שימוש ב-\en{Robust Cost Functions} המפחיתות
את השפעת מדידות חריגות.

\begin{table}[H]
\centering
\begin{tabular}{lcccc}
\toprule
מערכת & \en{ATE} [מ׳] & תדר [הרץ] & עומס מעבד [%] & חסינות לנשירה \\
\midrule
\en{LIO-SAM} & 0.18 & 15 & 50 & בינונית \\
\en{LVI-SAM} & 0.12 & 12 & 65 & גבוהה \\
\en{R2LIVE} & 0.08 & 20 & 55 & גבוהה \\
\en{NavigatorCrew} (מוצע) & 0.06 & 20 & 48 & גבוהה מאוד \\
\bottomrule
\end{tabular}
\caption{השוואת מערכות \en{Fusion} רב-חיישני על \en{M2DGR}.}
\label{tab:multi_sensor_fusion_comparison}
\end{table}

% ── 8.4 NavigatorCrew Fusion Architecture ───────────────────────────────────
\subsection{ארכיטקטורת ה-\en{Fusion} של \en{NavigatorCrew}}
\label{subsec:navcrew_fusion}

מערכת \en{NavigatorCrew} מממשת ארכיטקטורת \en{Fusion} היררכית
המשלבת את היתרונות של כל הגישות שתוארו. בשכבה התחתונה, כל
מודאליות חיישן מעובדת באופן עצמאי: \en{LOAM-AFAFC} ל-\en{LiDAR},
\en{ORB-SLAM3} למצלמה, ו-\en{IMU Preintegration} לאינרציאלי.
בשכבה העליונה, \en{Factor Graph} מבוסס \en{iSAM2} ממזג את
כל האילוצים לאופטימיזציה משותפת.

מנגנון ה-\en{Weighted Fusion} הדינמי מתאים את משקל כל חיישן
בהתאם לאיכות המדידה הנוכחית. לדוגמה, בסביבה חשוכה, משקל
המצלמה יופחת ומשקל ה-\en{LiDAR} יוגדל. בסביבה עם ערפל,
משקל ה-\en{LiDAR} יופחת ומשקל ה-\en{IMU} והמצלמה יוגדלו.

\begin{equation}
    \mathbf{X}^* = \arg\min_{\mathbf{X}} \sum_{m \in \mathcal{M}} w_m(t) \cdot \sum_{i \in \mathcal{I}_m} \| \mathbf{e}_{m,i}(\mathbf{X}) \|_{\Omega_{m,i}}^2
    \label{eq:navcrew_weighted_fusion}
\end{equation}

משוואה~\ref{eq:navcrew_weighted_fusion} מציגה את פונקציית
העלות המשוקללת, כאשר $\mathcal{M}$ היא קבוצת מודאליות החיישנים,
$w_m(t)$ הוא המשקל הדינמי של מודאליות $m$ בזמן $t$,
ו-$\mathcal{I}_m$ היא קבוצת המדידות ממודאליות $m$.

\begin{figure}[H]
    \centering
    \includegraphics[width=0.92\columnwidth]{figures/fig_navcrew_architecture.png}
    \caption{ארכיטקטורת ה-\en{Fusion} ההיררכית של \en{NavigatorCrew}.
             השכבה התחתונה מעבדת כל חיישן בנפרד; השכבה העליונה
             ממזגת את האילוצים באמצעות \en{Factor Graph}.}
    \label{fig:navcrew_architecture}
\end{figure}

% ── 8.5 Summary ─────────────────────────────────────────────────────────────
\subsection{סיכום}
\label{subsec:fusion_summary}

בפרק זה הצגנו את עקרונות ה-\en{Fusion} הרב-חיישני ואת המערכות
המובילות בתחום: \en{LIO-SAM}, \en{LVI-SAM} ו-\en{R2LIVE}.
מערכת \en{NavigatorCrew} המוצעת משלבת את היתרונות של גישות אלה
במסגרת \en{Fusion} היררכית עם משקלים דינמיים, ומשיגה דיוק
של 0.06 מטר על \en{M2DGR}. בפרק הבא נציג את ארכיטקטורת
המערכת המלאה ואת תוצאות הניסויים.